\chapter{JSON}
\label{ch:json}

\standfirst{Two containers, a parser that refuses more than it accepts, and
numbers that are still exact when they come out the other side.}

The \ct{Json} category reads and writes JSON: \ct{JSONObject} for a document's
names and values, \ct{JSONArray} for its lists, \ct{JSONParser} to read text
and \ct{JSONWriter} to write it. Everything raises \ct{JSONError} and nothing
else.

Most of the time you will not name any of those four. A String knows how to
read itself and every object knows how to write itself:

\begin{st}
'{"name": "Ada", "born": 1815}' asJson
\end{st}

\begin{st}
(Dictionary new at: 'k' put: 1; yourself) asJsonString
\end{st}

\ct{asJson} goes one way and \ct{asJsonString} the other, and they are
\emph{not} opposites---which is worth getting straight before anything else,
because both are defined on String. \ct{'[1,2]' asJson} is a \ct{JSONArray},
because there the String is a document being read. \ct{'[1,2]' asJsonString}
is the eight characters \ctp{"[1,2]"}, because there the String is a value
being written.

\section{Reading}

\ct{asJson} answers whatever the text is: an object, an array, a number, a
string, a boolean, or nil for \ct{null}. When you know which you want and
would rather be told than find out later, ask for it:

\begin{st}
JSONObject fromString: aString.
JSONArray fromString: aString.
JSONParser parse: aString
\end{st}

The first two raise if the document is not that kind. Reading a name is
\ct{at:}, and reading it with its type checked is one of seven accessors:

\begin{st}
| person |
person := '{"name": "Ada", "born": 1815, "tags": ["maths"]}' asJson.
person at: 'name'.                  "'Ada'"
person stringAt: 'name'.            "'Ada' -- and raises if it is not text"
person integerAt: 'born'.           "1815"
person arrayAt: 'tags'.             "a JSONArray"
(person arrayAt: 'tags') stringAt: 1
\end{st}

The seven are \ct{stringAt:}, \ct{numberAt:}, \ct{integerAt:}, \ct{floatAt:},
\ct{booleanAt:}, \ct{objectAt:} and \ct{arrayAt:}. Each has an
\ct{ifAbsent:} form, and each works the same way on a \ct{JSONArray} with an
index instead of a name.

\begin{stnote}[The wrong type raises]
\ct{stringAt:} of the number 42 does not answer \ct{'42'}. Java's JSON
libraries convert, and it is the kind of helpfulness that hides the sender's
bug in the reader's code: a document that says \ct{42} and one that says
\ct{"42"} are different documents, and the day one turns into the other is
the day to be told. The error names the key, what was found and what was
wanted.
\end{stnote}

\section{Names that are not there}

\ct{at:} answers nil for a name the document does not have. That is
deliberate and it is not what \ct{Dictionary} does: in JSON an absent field is
how \emph{optional} is spelled, and a reader that raised for one would turn
every read into \ct{at:ifAbsent:}.

The cost is that a missing name and a name whose value is \ct{null} both
answer nil. Two messages tell them apart and nothing else can:

\begin{st}
| document |
document := '{"middleName": null}' asJson.
document includesKey: 'middleName'.     "true"
document isNullAt: 'middleName'.        "true"
document includesKey: 'nickname'.       "false"
document isNullAt: 'nickname'           "false"
\end{st}

The difference is real. \ct|{"middleName": null}| says the sender knows about
middle names and this person has none. A document without the name says the
sender never considered the question.

\section{Writing}

\ct{asJsonString} answers compact text---no space anywhere it is not inside a
string---and \ct{asJsonStringIndent:} answers the indented form:

\begin{st}
aDocument asJsonString.
aDocument asJsonStringIndent: 2.
aDocument writeJsonOn: aStream
\end{st}

Both are on \ct{Object}, so nothing has to be built by hand first:

\begin{st}
#(1 2 3) asJsonString.                          "'[1,2,3]'"
(1 to: 3) asJsonString.                         "'[1,2,3]'"
'hello' asJsonString.                           "'""hello""'"
(Dictionary new at: 'k' put: 1; yourself) asJsonString
\end{st}

A \ct{Dictionary} becomes an object, every other collection becomes an array,
a \ct{Symbol} and a \ct{Character} become strings, and anything else raises.
A class of your own joins in by implementing \ct{asJsonValue} and answering
something that is already a JSON value.

\begin{stnote}[Names keep their order]
A \ct{JSONObject} writes its names in the order they were first put in. JSON
says the order is insignificant and \ct{diff} says otherwise: a document read
and written back comes out as it went in, and two runs of one program produce
the same file. 1983's \ct{Dictionary} iterates in hash order, so this is
something \ct{JSONObject} does rather than something it inherits.
\end{stnote}

\section{Numbers are exact}

A JSON number with a point or an exponent reads as a \ct{Fraction}, not a
\ct{Float}:

\begin{st}
'0.1' asJson.                   "1/10, and not 0.1"
'0.1' asJson * 10.              "1 exactly"
'1.5' asJson.                   "3/2"
'1.5' asJson asJsonString       "'1.5'"
\end{st}

This is the same decision \ct{DECIMAL} columns get in
Chapter~\ref{ch:database}, and JSON deserves it more, because JSON is what
money travels in now. A tenth read through a Float is already not a tenth;
the error surfaces as a penny three statements downstream rather than as a
failure where it happened.

A whole number answers an \ct{Integer} however it was written, because
Smalltalk's \ct{Fraction} reduces:

\begin{st}
'3.0' asJson.                   "3"
'1e2' asJson.                   "100"
'3.0' asJson isKindOf: Integer  "true"
\end{st}

And a JSON integer has no width, so an identifier from a database sequence
survives the trip:

\begin{st}
'123456789012345678901234567890' asJson
\end{st}

A C or a JavaScript reader would have lost its low digits, because both read
every number through a double.

When a Float is what you want, ask for one. \ct{floatAt:} is the lossy
accessor and it is spelled out at every call site, which is the point:

\begin{st}
('{"rate": 0.1}' asJson) floatAt: 'rate'
\end{st}

\begin{stwarn}
Writing a \ct{Float} back is not the same as writing what you read.
\ct{0.1 asJsonString} is the shortest decimal that Float rounds to, which is
\ctp{0.1}, and it is a different number from the one the document sent. Keep
the Fraction as long as you can and convert at the edge where the answer is
displayed.
\end{stwarn}

\section{What the parser refuses}

The grammar is RFC 8259 and nothing else. It is stricter than the JSON
libraries most people have met, and every one of these is something they
accept:

\begin{st}
'{"a":1,}' asJson.      "a trailing comma"
'{a:1}' asJson.         "an unquoted name"
'{''a'':1}' asJson.     "a single-quoted string"
'01' asJson.            "a leading zero"
'banana' asJson.        "a bare word"
'{"a":1} {"b":2}' asJson  "two documents where one was asked for"
\end{st}

All six raise. Leniency is not a kindness: it accepts a document that the
next program in the pipeline will reject, and the complaint then arrives from
somewhere that has never heard of this system. The sender of a malformed
document is the party who can fix it, and the way to tell them is to refuse
it.

Two limits are refusals of a different kind. A document may not nest deeper
than 512 levels, and an exponent may not exceed 4096---one so that forty
thousand open brackets cannot exhaust the heap, the other so that a
twelve-character document cannot ask this process for a billion digits. Both
matter because the document usually arrives from somewhere else. The first
can be raised deliberately:

\begin{st}
| parser |
parser := JSONParser on: aVeryDeepDocument.
parser maxDepth: 5000.
parser parse
\end{st}

\section{When it goes wrong}

Everything raises \ct{JSONError}, and every message from the parser carries
the character position:

\begin{st}
[aString asJson]
    on: JSONError
    do: [:error | Transcript show: error messageText; cr]
\end{st}

\begin{st}
'{"a":1,"b":}' asJson
\end{st}

\noindent answers, through the handler, \ctp{a value cannot begin with \$\} at
character 12}. A document that goes wrong at character 4096 of 40000 is one
somebody can open and look at; \ctp{invalid JSON} is not.

\section{Text is UTF-8}

A String in this system is bytes, so an escape has to become bytes.
\ct{\u0041} answers the one byte that is \ct{A}, and \ct{\u00e9} answers the
\emph{two} that are e-acute in UTF-8:

\begin{st}
'"\u00e9"' asJson size
\end{st}

\noindent is 2, not 1. A surrogate pair---which is how every emoji and a
great deal of Asian text arrives---is read as the one character it encodes,
and a lone surrogate is refused, because there is no encoding of one.

Going the other way, a byte above 126 is written as itself. Turning it back
into an escape would double the size of every document not written in English
to no purpose, since the two spellings are the same text to every reader.

\begin{stwarn}
\ct{\n} in a document is a line feed, 10, and it stays one. This image's text
panes break lines on the carriage return, 13, because that is what the Alto
did---so JSON text displayed straight in a Paragraph shows its line feeds as
control characters. Translating is the job of whatever puts the text on the
screen or in a file, and \ct{lib/Files-Fixes} already does it for files.
\end{stwarn}

\section{Building a document}

A \ct{JSONObject} is written into with \ct{at:put:}, which converts as it
goes---a \ct{Dictionary} becomes an object, an \ct{Array} becomes an array,
and anything with no JSON spelling is refused right there:

\begin{st}
| reply |
reply := JSONObject new.
reply at: 'ok' put: true.
reply at: 'rows' put: #(1 2 3).
reply at: 'error' put: nil.
reply asJsonString
\end{st}

Refusing at the put rather than at the write is the whole reason the
conversion is a message. In a web application those two moments are a request
apart, and only the first one has the code that was wrong on the stack.

A \ct{JSONArray} takes \ct{add:}, \ct{addAll:} and \ct{at:put:}, and answers
to enough of the collection protocol to be useful:

\begin{st}
| numbers |
numbers := '[1,2,3,4]' asJson.
(numbers select: [:each | each odd]) asJsonString.    "'[1,3]'"
numbers inject: 0 into: [:a :b | a + b].              "10"
numbers asArray
\end{st}

\begin{stnote}[Two places it is not a collection]
\ct{collect:} answers an \ct{OrderedCollection} and not a \ct{JSONArray},
because the block may answer anything at all and a \ct{JSONArray} may hold
only JSON values---\ct{collect: [:each | each class]} would otherwise raise
on a perfectly reasonable block. And an index past the end raises rather than
answering nil: an array has a length, and reading past it is a mistake in
your arithmetic rather than an optional field.
\end{stnote}

\section{More than one process at once}

A \ct{JSONObject} and a \ct{JSONArray} are safe to use from several processes
at once, which on this system means several cores. That is not what the base
collections do---Chapter~\ref{ch:contract} explains why an
\ct{OrderedCollection} is deliberately unsynchronized---and it is a decision
made for this one case: a JSON document is the object a request arrives as and
is then handed to whatever answers it, so it is shared by nature rather than by
accident.

What you get is that no message loses anything. Two processes putting two names
both keep them; two processes adding to one array both keep their elements; a
process reading while another writes sees the value before or the value after.

What you do not get is a sequence:

\begin{st}
document at: 'n' put: (document at: 'n') + 1
\end{st}

\noindent is a read and a write with a gap, and two processes doing it lose an
increment---exactly as they would with any other object. Hold your own lock
around a sequence that has to happen together. \ct{at:ifAbsentPut:} is provided
because that particular sequence is common enough to be worth doing properly.

\begin{stnote}[Your block never runs inside the lock]
\ct{do:}, \ct{keysAndValuesDo:}, \ct{collect:} and \ct{select:} take a
snapshot and then enumerate it, so your block sees the members as they were when
it started---and so that \ct{json do: [:each | json at: ...]} works. A
\ct{Mutex} here is not re-entrant, so a design that held the lock across your
block would have made the most natural line somebody could write into a
deadlock.
\end{stnote}

\begin{stnote}
An object with thousands of \emph{names} is linear to build, because
\ct{lib/Collections-Protocol} gives \ct{String} a hash over every character.
1983's \ct{String>>hash} reads the first character, the second-to-last and the
length, so two hundred names like \ct{key1}..\ct{key200} share eleven hash
values and a \ct{Dictionary} of them degenerates into one long chain. See
\ct{doc/JSON.md}.
\end{stnote}

\section{Streams}

The parser and the writer both take a stream, so a document does not have to
be a String first:

\begin{st}
JSONParser parse: aReadStream.
aValue writeJsonOn: aWriteStream
\end{st}

A stream can also be read one document at a time, which is what a log file of
them is:

\begin{st}
| parser |
parser := JSONParser on: (ReadStream on: '[1] [2]').
parser nextValue.       "the first array"
parser nextValue        "the second"
\end{st}

\ct{parse} reads one value and insists the text ends there. \ct{nextValue}
reads one value and leaves the stream where it stopped.

\section{The whole protocol}

\begin{center}
\begin{tabular}{@{}ll@{}}
\toprule
\textbf{Reading} & \\
\ctp{aString asJson} & the value that text spells \\
\ctp{JSONObject fromString:} & and it must be an object \\
\ctp{JSONArray fromString:} & and it must be an array \\
\ctp{JSONParser parse:} & a String or a stream \\
\ctp{JSONParser on:} & a parser, for maxDepth: and nextValue \\
\midrule
\textbf{Writing} & \\
\ctp{anObject asJsonString} & compact \\
\ctp{anObject asJsonStringIndent:} & one member to a line \\
\ctp{anObject writeJsonOn:} & straight onto a stream \\
\ctp{anObject asJsonValue} & the conversion, for your own classes \\
\midrule
\textbf{Both containers} & \\
\ctp{at:} \ctp{at:ifAbsent:} \ctp{at:put:} & by name, or by index \\
\ctp{size} \ctp{isEmpty} \ctp{notEmpty} \ctp{do:} & \\
\ctp{stringAt:} \ctp{numberAt:} \ctp{integerAt:} & each with an ifAbsent: form \\
\ctp{floatAt:} \ctp{booleanAt:} & \\
\ctp{objectAt:} \ctp{arrayAt:} & \\
\midrule
\textbf{JSONObject} & \\
\ctp{includesKey:} \ctp{isNullAt:} & absent, or present and null \\
\ctp{keys} \ctp{values} \ctp{keysAndValuesDo:} & in the order they were put \\
\ctp{removeKey:} \ctp{at:ifAbsentPut:} \ctp{asDictionary} & \\
\midrule
\textbf{JSONArray} & \\
\ctp{add:} \ctp{addAll:} \ctp{first} \ctp{last} & \\
\ctp{removeAt:} \ctp{removeFirst} \ctp{removeLast} & \\
\ctp{select:} \ctp{reject:} \ctp{collect:} \ctp{detect:ifNone:} & collect: answers an OrderedCollection \\
\ctp{inject:into:} \ctp{asArray} \ctp{asOrderedCollection} & \\
\bottomrule
\end{tabular}
\end{center}

\section{Where it came from, and what it is not}

The API is the shape \ct{org.kissweb.json} has, which is the shape every JSON
library has had since 2002. None of the code is: that package is a fork of
JSON-java and carries a licence adding \emph{The Software shall be used for
Good, not Evil} to MIT, which is why it is not free software and why not one
line of it could be taken. \ct{lib/Json/PROVENANCE.md} records what was kept
of the shape, what was deliberately dropped, and why.

What this package is not is a JSON \emph{pointer}, a JSON \emph{path}, or a
converter to and from XML. Those exist in the Java and are not JSON; if they
are ever wanted here they belong in packages of their own.
